\documentclass[a4paper,12pt]{article}

\usepackage[utf8]{inputenc}
\usepackage[spanish]{babel}
\usepackage{exerlist}
\usepackage{math2}
\usepackage{amsmath}

\newcommand{\myneg}{\raise.17ex\hbox{$\scriptstyle\sim$}}%
\newcommand{\mysmallneg}{\raise.17ex\hbox{$\scriptscriptstyle\sim$}}%
\newcommand{\codlength}[1]{\left< #1 \right>}
\newcommand{\colvector}[2]{{#1 \choose #2}}
\newcommand{\lineal}[1]{\mbox{lin}(#1)}
\newcommand{\Ppolar}[1]{#1^{\bigtriangleup}}
\renewcommand{\conv}[1]{\mathrm{conv}\left(#1\right)}
\renewcommand{\cone}[1]{\mathrm{cone}\left(#1\right)}
\renewcommand{\lin}[1]{\mathrm{lin}\left(#1\right)}
\newcommand{\homog}[1]{\mathrm{homog}\left(#1\right)}
\newcommand{\rec}[1]{\mathrm{rec}\left(#1\right)}
\newcommand{\HH}{\mathcal{H}}
\newcommand{\VV}{\mathcal{V}}
\newcommand{\proj}{\mathrm{proj}}

\setlecture{Geometría Discreta}
\setprofesor{Luis M. Torres}
\setexamdate{Julio, 2025}

\begin{document}
\exercisetitle{Ejercicios Capítulo 2}

\begin{exerciselist}

  \item Sean $Y = \left( \begin{array}{cc} 1 & 0 \\ 1 & 1 \end{array}\right)$ y $C$ el
  $\mathcal{V}$-cono definido por $C := \mathrm{cone}(Y)$. Determinar un sistema de
  desigualdades lineales en las variables $x_1, x_2, t_1, t_2$, asociado a un $\mathcal{H}$-cono $\hat{C} \subset \RR^4$
  con la propiedad de que $C \cong \mathrm{proj}_{t_1}(\mathrm{proj}_{t_2}(\hat{C}))$.
  Aplicar el método de Fourier-Motzkin para eliminar las variables $t_1$ y $t_2$.
  de este sistema.

  \item Sea $P := \conv{V} + \cone{Y} \subset \RR^2$ un $\VV$-poliedro, con
  $$
  V = \left(
  \begin{array}{cc}
    0 & 0 \\
    1  & -1
  \end{array}
  \right),
  \qquad
  Y = \left(
  \begin{array}{c}
    1 \\
    0
  \end{array}
  \right),
  \qquad
  $$
  \begin{exerciselist}
    \item Dibujar $P$.
    \item Expresar el cono de homogeneización $C:=\homog{P}$ en su forma $\VV$.
    \item Construir un $\HH$-cono $\hat{C} \subset \RR^6$, dado en las variables $x_0, x_1, x_2, t_1, t_2, t_3$ con la propiedad de que $C \cong \proj_{t_1}(\proj_{t_2}(\proj_{t_3}(\hat{C})))$.
    \item Utilizar el método de Fourier-Motzkin para expresar $C$ en la forma $\HH$.
    \item Expresar $P$ en la forma $\HH$.
  \end{exerciselist}

  \item Sean
  $$
  V:= \left(
  \begin{array}{cc}
  1 & 0 \\
  0 & 1
  \end{array}
  \right),
  \qquad
  Y:= \colvector{1}{3},
  \qquad
  P:= \conv{V} + \cone{Y}.
  $$

  \begin{exerciselist}
  \item Encontrar un $\mathcal{H}$-poliedro $\hat{P}:= \left\{ w:= (x_1, x_2,
  \lambda_1, \lambda_2, t)^T \in \RR^5 \, : \, \hat{A} w \leq \hat{b} \right\}$,
  con $\hat{A} \in \RR^{M \times 5}$ y $\hat{b} \in \RR^M$ tal que $P$ sea
  congruente con la proyección de $\hat{P}$ en la dirección de los ejes
  coordenados correspondientes a las variables $\lambda_1, \lambda_2$ y $t$.
  \item Utilizar el algoritmo de elimnación de Fourier-Motzkin para escribir a
  $P$ en la forma $\mathcal{H}$.
  \end{exerciselist}

  \item Sea $P \subset \RR^2$ el conjunto solución del siguiente sistema de desigualdades lineales:
  $$
  \left\{
  \begin{array}{l}
  x_1 + x_2 \leq 1,\\
  - x_1 + x_2 \leq 1.
  \end{array}
  \right.
  $$
  \begin{exerciselist}
    \item Dibujar $P$.
    \item Expresar el cono de homogeneización $C:=\homog{P}$ en su forma $\HH$.
    \item Construir un $\VV$-cono $\hat{C} \subset \RR^5$, dado en las variables $x_0, x_1, x_2, w_1, w_2$ 
    con la propiedad de que $C \cong \hat{C} \cap H_{w_1} \cap H_{w_2}$.
    \item Expresar $C$ en la forma $\VV$.
    \item Expresar $P$ en la forma $\VV$.
  \end{exerciselist}

  \item
  Sea $C \subset \RR^3$ el conjunto solución del siguiente sistema de desigualdades lineales.
  $$
  \left\{
  \begin{array}{l}
  - x_1 + x_2 - x_3 \leq 0,\\
  - x_1 - x_2 + x_3 \leq 0,\\
   \,\,\,\,\, x_1 - x_2 - x_3 \leq 0
  \end{array}
  \right.
  $$

  \begin{exerciselist}
  \item Encontrar un conjunto finito de vectores $\hat{Y} \subset \RR^6$ tal que el
  $\mathcal{V}$-cono generado por ellos $\hat{C}:= \left\{ (x_1, x_2, x_3, w_1, w_2, w_3)^T \in \cone{\hat{Y}} \right\}$ satisfaga la propiedad:
  $$
  C \cong \hat{C} \cap \left( \bigcap_{i=1}^3 \{ w_i = 0 \} \right).
  $$
  \item Calcular los vectores generadores de $\hat{C} \cap \{w_3 = 0 \}$.
  \end{exerciselist}

  \item Sean
  $$
  V:= \left(
  \begin{array}{rrrrrr}
  40  & 24  & 12 & -12 & -30 & 2\\
  -30 & -24 & -6 &   12 &  30 & -1\\
  3    & 2    &  1 &  0  &  -1   & 0
  \end{array}
  \right), \qquad
  P:= \conv{V} \subset \RR^3.
  $$
  Calcular un conjunto (matriz) $V'$ tal que
  $$
  P \cap \setof{x \in \RR^3}{x_3 = 0} = \conv{V'}.
  $$
  Dibujar la región $\conv{V'}$.

  \item Sean
  $$
  V:= \left(
  \begin{array}{ccc}
  0 & 1 & 1 \\
  1 & 1 & -1
  \end{array}
  \right), \qquad
  P:= \conv{V} \subset \RR^2.
  $$
  Determinar un poliedro $\hat{P}=\setof{x \in \RR^d}{\hat{A}x \leq \hat{b}}$,
  con $d>2$, que satisfaga la propiedad de que $P$ sea congruente con el
  poliedro obtenido a partir de proyectar $\hat{P}$ sucesivamente en la
  dirección de algunos ejes coordenados. Escribir el sistema $\hat{A}x \leq
  \hat{b}$.

  \item Demostrar la siguiente versión del Lema de Farkas: Dados $A \in \RR^{m \times d}$ y $b \in \RR^m$,
  exactamente uno de los dos enunciados siguientes es verdadero:
  \begin{enumerate}
  \item Existe $x \in \RR^d$ tal que $Ax=b$, $x \geq \zero$.
  \item Existe $y \in \RR^m$ tal que
  $y^TA \geq \zero^T$, $y^Tb < 0$.
  \end{enumerate}

  \item Demostrar la siguiente versión del Lema de Farkas: Dados $A \in \RR^{m_1
  \times d}, C \in \RR^{m_2 \times d}$, $b \in \RR^{m_1}$ y $f \in \RR^{m_2}$,
  exactamente uno de los dos enunciados siguientes es verdadero:
  \begin{enumerate}
  \item Existe $x \in \RR^d$ tal que $Ax=b$, $Cx \geq f$.
  \item Existen $y \in \RR^{m_1}$, $w \in \RR^{m_2}$, $w \leq 0$ tales que
  $y^TA + w^TC = \mathbf{0}^T$, $y^Tb + w^Tf < 0$.
  \end{enumerate}

  \item Formular y demostrar una versión del Lema de Farkas que caracterice
  cuándo un sistema de la forma $Ax \leq b$, $x \geq 0$ tiene solución.

  \item Usar el Lema de Farkas para demostrar el siguiente resultado.

  Dada una matriz $A \in \RR^{m \times d}$, uno y sólo uno de las dos enunciados siguientes es verdadero:
  \begin{enumerate}
  \item Existe $x \in \RR^d$ tal que $Ax < 0$.
  \item Existe $y \in \RR^m$, $y \geq 0$, $y \neq 0$ tal que $y^T A = 0$.
  \end{enumerate}

  \textbf{Nota:} Si $w$ es un vector, usamos la notación $w < 0$ para decir que \emph{cada una de las componentes} de $w$ es estrictamente menor que cero. Observar que esto es equivalente a afirmar que existe al menos un $\varepsilon > 0$  tal que $w \leq - \varepsilon \mathbf{1}$.

  \item Sea $P:=P(A,b) \subset \RR^d$ un $\HH$-poliedro, con $A \in \RR^{m \times d}$ y $b \in \RR^m$. Definimos
  $$
  C(P):= P\left(
  \left(
  \begin{array}{cc}
    -1 & \zero^T \\
    -b & A \\
  \end{array}
  \right),
  \colvector{0}{\zero}
  \right) \subset \RR^{d+1}.
  $$
  Demostrar que:
  \begin{exerciselist}
  \item $$
  P = \setof{x \in \RR^d}{\colvector{1}{x} \in C(P)}.
  $$
  \item $C(P) \subset H_0^{+}$, donde $H_0^{+} := \setof{\colvector{x0}{x} \in \RR^{d+1}}{x_0 \geq 0}$.
  \end{exerciselist}
  

  \item Sea $\hat{P}:=P(\hat{A},\hat{b}) \subset \RR^{d+1}$ un $\HH$-poliedro, con $\hat{A} \in \RR^{m \times (d+1)}$ y $\hat{b} \in \RR^m$. Demostrar que
  $$
  P := \setof{x \in \RR^d}{\colvector{1}{x} \in \hat{P}}
  $$
  es un $\HH$-poliedro en $\RR^d$.

  \item Sea $P:=\conv{V} + \cone{Y} \subset \RR^d$ un $\VV$-poliedro, con $V \in \RR^{d \times n}$ y $Y \in \RR^{d \times k}$. Definimos
  $$
  C(P):= \cone{
  \begin{array}{cc}
    \one^T & \zero^T \\
    V & Y
  \end{array}
  } \subset \RR^{d+1}.
  $$
  Demostrar que:
  $$
  P = \setof{x \in \RR^d}{\colvector{1}{x} \in C(P)}.
  $$

  \item Sea $P:=\conv{V} \subset \RR^d$ un $\VV$-polítopo, con $V \in \RR^{d \times n}$. Definimos
  $$
  C(P):= \cone{
  \begin{array}{cc}
    \one^T \\ V
  \end{array}
  } \subset \RR^{d+1}.
  $$
  Demostrar que:
  $$
  P = \setof{x \in \RR^d}{\colvector{1}{x} \in C(P)}.
  $$

  \item Sean $W \in \RR^{(d+1) \times k}$ y $\hat{C} = \mathrm{cone}(W) \subset
  H_0^{+}$, donde $H_0^{+} := \setof{\colvector{x0}{x} \in \RR^{d+1}}{x_0 \geq 0}$. 
  Demostrar que el conjunto:
  $$
  P := \setof{x \in \RR^d}{{1 \choose x} \in \hat{C}}
  $$
  es un $\mathcal{V}$-poliedro en $\RR^d$.

  \item Sea $P:= P(A, b) \subset \RR^d$ un $\HH$-poliedro, con $A \in \RR^{m \times d}$ y $b \in \RR^m$. Definimos el $\HH$-cono
  $$
  Q := P((-b \,\,\,\, A), \zero) \subset \RR^{d+1}.
  $$
  Demostrar que:
  $$
  P = \emptyset
  \qquad \Leftrightarrow \qquad
  Q \subseteq \setof{x \in \RR^{d+1}}{x_0 \leq 0},
  $$
  donde $x_0$ designa a la primera componente de $x \in \RR^{d+1}$.

  \item Usando los resultados de los ejercicios anteriores, demostrar que el resultado del Teorema de Minkowski-Weyl para poliedros se obtiene como corolario del Teorema de Minkowski-Weyl para conos.

  \item Demostrar que el Teorema de Minkowski-Weyl para polítopos se obtiene como corolario del Teorema de Minkowski-Weyl para poliedros.

  \item Sea $P:=P(A,b) \subset \RR^d$ un $\HH$-poliedro, con $A \in \RR^{m \times d}$ y $b \in \RR^m$. Demostrar que $\lin{P} = \setof{x \in R^d}{Ax=0}$, es decir, que el espacio de linealidad de $P$ es el espacio núcleo de $A$.

  \item Sea $P \subset \RR^d$ un poliedro. Demostrar que si $U$ es el complemento ortogonal de $\lin{P}$, entonces se cumple que $\lin{P \cap U} = \{ \zero \}$.

  \item Sea $P \subset \RR^d$ un poliedro. Su cono de homogeneización está definido por $\homog{P}:= C_1 + C_2$, con
  $$
  C_1:= \setof{\alpha {1 \choose x} \in \RR^{d+1}}{x \in P, \alpha \geq 0},
  \qquad
  C_2:= \setof{{0 \choose y} \in \RR^{d+1}}{y \in \rec{P}}.
  $$
  \begin{exerciselist}
    \item Demostrar que $C_1$ es un cono convexo.
    \item Demostrar que $C_2$ es un cono convexo. (Puede darse por conocido que
    $\rec{P}$ es un cono convexo).
    \item Demostrar que la suma de Minkowski de dos conos convexos es un cono
    convexo, de donde se concluye que $\homog{P}$ es un cono convexo.
  \end{exerciselist}

 \item Sea $P \subset \RR^d$ un conjunto convexo.
  \begin{exerciselist}
    \item Demostrar que $\rec{P}$ es un cono convexo.
    \item Demostrar que $\homog{P}$ es un cono convexo en $\RR^{d+1}$.
    \item Demostrar que
    $$
    P = \setof{x \in \RR^d}{\colvector{1}{x} \in \homog{P}}.
    $$
  \end{exerciselist}

\item Sea $P:=\conv{V} + \cone{Y} \subset \RR^d$ un $\VV$-poliedro, con $V \in \RR^{d \times n}$ y $Y \in \RR^{d \times k}$. Demostrar que:
$$
\homog{P}= \cone{
\begin{array}{cc}
  \one^T & \zero^T \\
  V & Y
\end{array}
}.
$$

\item Empleando el Teorema de Carathéodory para conos, demostrar como corolario el Teorema de Carathéodory para polítopos:
\begin{quote}
  Sean $V \in \RR^{d \times n}$ y $x \in \RR^d$. Si $x \in \conv{V}$, entonces existe $V' \subset V$, con $\card{V'} \leq \dim(\conv{V}) + 1 = \mathrm{rg}{\one^T \choose V}$ tal que $x \in \conv{V'}$.
\end{quote}
\end{exerciselist}

\end{document}
